\section{Sign conventions for the Hamiltonian comparison}
\label{sec:app-sign-conventions}
\label{app:sign-conventions}

The two \(\mathfrak h\)-modules attached to the labels \((a,b)\) go in
opposite directions.  The polynomial PBW special-fibre descendant
\(\widetilde\Psi_{a,b}\), symmetrised trace projection of the open
Hamiltonian action on the matrix Koszul algebra \(A_\psi\), admits the
lowering action
\(z_1\cdot\widetilde\Psi_{a,b}=b\,\widetilde\Psi_{a,b-1}\),
\(z_2\cdot\widetilde\Psi_{a,b}=-a\,\widetilde\Psi_{a-1,b}\).  The
Matlis principal-part vector
\(\rho_{a,b}=z_1^{-a-1}z_2^{-b-1}dz_1dz_2\), residue dual of the
formal Hamiltonian polydisk against its continuous Taylor dual, admits
the raising action \(z_1\cdot\rho_{a,b}=-(b+1)\rho_{a,b+1}\),
\(z_2\cdot\rho_{a,b}=(a+1)\rho_{a+1,b}\).  Diagonal rescaling
\(\rho_{a,b}\mapsto\lambda_{a,b}\rho_{a,b}\),
\(\widetilde\Psi_{a,b}\mapsto\mu_{a,b}\widetilde\Psi_{a,b}\) preserves
the locally nilpotent lowering direction on descendants and the raising
direction on principal parts; the directions of the index shifts already
disagree.

The cohomological shift convention is \((V[1])^n=V^{n+1}\); the
reduced matrix Koszul letter \(\psi\) is odd, and cyclic graded signs
use its parity.  Absolute BV grading enters only through the global
shift of the cotangent coordinate
(Definition~\ref{def:app-unreduced-bv-degrees}).  BV degree, ghost
number, and form degree are tracked separately: BV degree is
cohomological, ghost number is the antifield/ghost grading of the
classical fields \((\phi_1,\phi_2,\psi)\) and
\((\mathcal A,\mathcal B)\), and form degree is the de Rham/Dolbeault
exterior grading on \(\Omega^\bullet(\R^2)\widehat\otimes
\Omega^{0,\bullet}(\C^2)\); none of the three is identified with
another.  Hochschild homology \(HH_\bullet\), cyclic homology, and
negative-cyclic homology \(HC^-_\bullet\) are distinguished at every
reference; the cyclic trace symbol \(\operatorname{Tr}\) below is the
graded cyclic quotient of
Definition~\ref{def:app-sign-hamiltonian-cyclic}.  Its Hochschild
class is obtained by forgetting the \(S^1\)-equivariance.  Connes'
\(B\) is the degree-\(+1\) operator in the mixed complex and the
connecting operator in the SBI sequence.  A negative-cyclic class is
the constellation-level Calabi--Yau pairing only after a compact
CY-to-chiral target category and its negative-cyclic trace class have
been supplied.  The formal Hamiltonian stalk uses the graded cyclic
trace quotient and its underlying Hochschild class.  The BCOV holomorphic-anomaly
equation enters only as a
target equation on a compact Calabi--Yau threefold, with the
BCOV-strict sign of \cite{bcov} preserved by the matched-conventions
divergence row of Theorem~\ref{thm:divergence-bcov}; the local
Hamiltonian theorem is the formal-Darboux stalk.

\begin{defn}[Local symplectic and cyclic conventions]
\label{def:app-sign-hamiltonian-cyclic}
  Fix the formal-local holomorphic symplectic datum
  \(\bigl(\widehat{\C^2}_0,\Omega_{\mathrm{loc}}=dz_1\wedge dz_2\bigr)\)
  of Section~\ref{sec:app-matlis-principal-parts}.  The Hamiltonian
  vector field of \(f\in\C[[z_1,z_2]]\) is determined by
  \[
    \iota_{X_f}\Omega_{\mathrm{loc}}=-df,
  \]
  giving
  \[
    X_f=(\partial_{z_1}f)\partial_{z_2}
        -(\partial_{z_2}f)\partial_{z_1},
    \qquad
    X_{z_1}=\partial_{z_2},\quad X_{z_2}=-\partial_{z_1}.
  \]
  The Poisson bracket is
  \[
    \{f,g\}
      =X_f(g)
      =\partial_{z_1}f\,\partial_{z_2}g
      -\partial_{z_2}f\,\partial_{z_1}g,
  \]
  in particular \(\{z_1,z_2\}=1\).  Hamiltonian fields are
  divergence-free for \(\Omega_{\mathrm{loc}}\): the volume-form scalar
  contribution is zero
  (Proposition~\ref{prop:app-local-geometry-scalar-shadow}).

  The reduced matrix Koszul algebra of
  Definition~\ref{def:app-full-psi-complex} is
  \[
    A_\psi=\C\langle x,y,\psi\rangle,
    \qquad
    |x|=|y|=0,\quad |\psi|\equiv1\pmod 2,
    \qquad
    Q\psi=xy-yx,\ \ Qx=Qy=0.
  \]
  The graded cyclic quotient uses the Koszul sign rule
  \[
    [uv]_{\mathrm{cyc}}=(-1)^{|u||v|}[vu]_{\mathrm{cyc}},
  \]
  and the trace symbol \(\operatorname{Tr}\) is graded cyclic with the
  same sign.  The only odd letter is \(\psi\).
\end{defn}

\begin{lem}[Chevalley--Eilenberg and Lichnerowicz--Poisson gradings]
\label{lem:app-sign-ce-pv}
  Let \(\mathfrak l\) be a finite-dimensional Lie algebra with basis
  \(e_I\), dual basis \(\epsilon^I\), and structure constants
  \([e_I,e_J]=C^K_{IJ}e_K\).  For
  \(T^*[1]\mathfrak l=\mathfrak l\ltimes\mathfrak l^\vee[1]\), the
  standard Chevalley--Eilenberg algebra has generators
  \[
    c^I=e_I^\vee[-1],\qquad u_I=(\epsilon^I[1])^\vee[-1],
    \qquad |c^I|_{\mathrm{CE}}=1,\quad |u_I|_{\mathrm{CE}}=2.
  \]
  Give \(c^I\) cotangent weight \(0\) and \(u_I\) cotangent weight \(1\),
  and define
  \[
    \operatorname{Regr}_u C^\bullet_{\mathrm{CE}}(T^*[1]\mathfrak l)
    \quad\text{by}\quad
    |a|_{\mathrm{PV}}=|a|_{\mathrm{CE}}-2\operatorname{wt}_u(a).
  \]
  The differential preserves cotangent weight and is
  \[
    d_{\mathrm{CE}}c^K=-\tfrac12 C^K_{IJ}c^Ic^J,\qquad
    d_{\mathrm{CE}}u_K=C^L_{KJ}u_Lc^J.
  \]
  Let \(A_{\mathfrak l}=\mathbf S(\mathfrak l)\).  Its
  Lichnerowicz--Poisson complex has even linear functions \(O_K\), odd
  constant vector fields \(\partial^I\), Schouten bracket
  \([\partial^I,O_J]_S=\delta^I_J\), and linear Poisson bivector
  \[
    \pi=\tfrac12 C^K_{IJ}O_K\partial^I\partial^J.
  \]
  Then
  \[
    d_\pi\partial^K=-\tfrac12 C^K_{IJ}\partial^I\partial^J,
    \qquad
    d_\pi O_K=C^L_{KJ}O_L\partial^J,
  \]
  and the assignment
  \[
    c^I\longmapsto\partial^I,\qquad u_K\longmapsto O_K
  \]
  is an isomorphism of regraded cochain algebras
  \[
    \operatorname{Regr}_u C^\bullet_{\mathrm{CE}}(T^*[1]\mathfrak l)
      \xrightarrow{\sim}
    \bigl(\mathrm{PV}(A_{\mathfrak l}),d_\pi\bigr).
  \]
  The Hamiltonian vector field of \(O_K\) is
  \(X_K=d_\pi O_K\); it is a linear combination of the constant fields
  \(\partial^J\), rather than a generator \(\partial^K\).
\end{lem}

\begin{proof}
  The shift convention \((V[1])^n=V^{n+1}\) places
  \(\mathfrak l^\vee[1]\) in degree \(-1\), hence its standard
  Chevalley--Eilenberg dual coordinate has degree \(2\).  The two
  differential formulas follow from the Lie bracket and the right
  coadjoint action.  They preserve cotangent weight, so the regraded
  differential still has degree \(+1\).

  The Schouten identities follow from
  \([\pi,O_K]_S=\partial_{\partial^K}\pi\) and
  \([\pi,\partial^K]_S=-\partial_{O_K}\pi\).  Thus the displayed
  substitution intertwines the differentials on generators.  Both
  differentials are derivations, which proves the cochain-algebra
  isomorphism.  Finally,
  \(d_\pi O_K=C^L_{KJ}O_L\partial^J\); for an abelian
  \(\mathfrak l\), this Hamiltonian field is zero while every
  \(\partial^K\) remains a nonzero constant vector field.  Hence the two
  kinds of vector field are intrinsically distinct.
\end{proof}

\begin{lem}[Primitive one-\(\psi\) chain boundary]
\label{lem:app-sign-one-psi-boundary}
  In \(A_\psi\) of
  Definition~\ref{def:app-sign-hamiltonian-cyclic}, for ordinary words
  \(u,v\) in \(x,y\),
  \[
    Q\,\operatorname{Tr}(\psi u\psi v)
      =
      \operatorname{Tr}\bigl(\psi vxyu-\psi vyxu-\psi uxyv+\psi uyxv\bigr).
  \]
  Removing the common front \(\psi\) gives the boundary in the
  one-\(\psi\) chain model
  \[
    \partial_2(u,v)=vxyu-vyxu-uxyv+uyxv.
  \]
\end{lem}

\begin{proof}
  The first displayed \(\psi\) contributes \(+(xy-yx)\) under \(Q\);
  the second contributes \(-(xy-yx)\) because exactly one odd letter
  precedes it, so the Koszul sign of the derivation flips.  The even
  word produced by the first replacement may be moved cyclically past
  the remaining \(\psi\) and the word \(v\) with sign \(+1\):
  the Koszul sign \((-1)^{|u||v|}\) for moving an even string past a
  graded cyclic block is trivial.  Removing the common front \(\psi\)
  identifies the result with the displayed boundary.  Cross-checked
  against Lemma~\ref{lem:app-full-two-psi-boundary-sign}.
\end{proof}

\begin{lem}[Residue coadjoint signs]
\label{lem:app-sign-principal-part}
  Let
  \[
    \rho_{a,b}=z_1^{-a-1}z_2^{-b-1}dz_1dz_2,\qquad a,b\geq0,
  \]
  with residue pairing
  \(\langle\rho_{a,b},z_1^mz_2^n\rangle=\delta_{a,m}\delta_{b,n}\) and
  coadjoint sign \((f\cdot\rho)(g)=-\rho(\{f,g\})\).  Then
  \[
    z_1^pz_2^q\cdot\rho_{a,b}
      =-(pb-qa+p-q)\rho_{a-p+1,b-q+1},
  \]
  with negative-index targets and the scalar target \(\rho_{0,0}\)
  set to zero.  In particular,
  \[
    z_1\cdot\rho_{a,b}=-(b+1)\rho_{a,b+1},\qquad
    z_2\cdot\rho_{a,b}=(a+1)\rho_{a+1,b}.
  \]
\end{lem}

\begin{proof}
  Divergence-free Hamiltonian fields
  (Definition~\ref{def:app-sign-hamiltonian-cyclic}) make the residue
  pairing \(\mathfrak h\)-equivariant via integration by parts with zero
  volume-form correction, fixing the coadjoint convention
  \((f\cdot\rho)(g)=-\rho(\{f,g\})\).  Pair against a test
  monomial \(g=z_1^mz_2^n\):
  \[
    (z_1^pz_2^q\cdot\rho_{a,b})(z_1^mz_2^n)
      =-\rho_{a,b}(\{z_1^pz_2^q,z_1^mz_2^n\}).
  \]
  Direct computation gives
  \(\{z_1^pz_2^q,z_1^mz_2^n\}=(pn-qm)z_1^{p+m-1}z_2^{q+n-1}\), and the
  residue pairing keeps the \((a,b)\) Taylor coefficient.  The unique
  \((m,n)\) producing \(z_1^az_2^b\) is \(m=a-p+1,\ n=b-q+1\);
  substitution yields coefficient
  \(p(b-q+1)-q(a-p+1)=pb-qa+p-q\), and the leading minus from the
  coadjoint convention gives the displayed formula.  Targets with
  \(a-p+1<0\) or \(b-q+1<0\) require \(m<0\) or \(n<0\) and are
  unreachable by polynomial test monomials, hence the cutoff at
  negative indices.  The scalar target \(\rho_{0,0}\) is the case
  \(p=a+1,\ q=b+1\), at which the coefficient
  \((a+1)b-(b+1)a+(a+1)-(b+1)=0\) vanishes automatically; the formula
  is therefore consistent with the exclusion of \(\rho_{0,0}\) from
  \(\mathcal P=\operatorname{Ann}_{\mathrm{res}}(\C\cdot1)\).
\end{proof}

\begin{lem}[PBW descendant signs]
\label{lem:app-sign-pbw-descendant}
  The polynomial primitive one-\(\psi\) descendants
  \(\widetilde\Psi_{a,b}\) of
  Corollary~\ref{cor:descendant-coadjoint-difference} carry the
  symmetric PBW special-fibre action
  \[
    z_1^pz_2^q:\widetilde\Psi_{a,b}
      \longmapsto
      (pb-qa)\widetilde\Psi_{a+p-1,b+q-1}.
  \]
  In particular,
  \(z_1\cdot\widetilde\Psi_{a,b}=b\widetilde\Psi_{a,b-1}\) and
  \(z_2\cdot\widetilde\Psi_{a,b}=-a\widetilde\Psi_{a-1,b}\): the action
  lowers the polynomial bidegree.
\end{lem}

\begin{proof}
  The Hamiltonian vector field of \(z_1^pz_2^q\) is
  \[
    X_{z_1^pz_2^q}=pz_1^{p-1}z_2^q\partial_{z_2}
                  -qz_1^pz_2^{q-1}\partial_{z_1}.
  \]
  Acting on the symmetrised trace with \(a\) factors of \(x\) and
  \(b\) factors of \(y\), the first summand differentiates one of the
  \(b\) copies of \(y\) and contributes \(pb\) copies of the
  symmetrised target \(\widetilde\Psi_{a+p-1,b+q-1}\); the second
  differentiates one of the \(a\) copies of \(x\) and contributes
  \(-qa\) copies of the same target.  The sum is the displayed
  coefficient.
\end{proof}

\begin{lem}[Moyal expansion coefficients]
\label{lem:app-sign-moyal-coefficients}
  With
  \[
    P=\partial_{z_1}\otimes\partial_{z_2}
      -\partial_{z_2}\otimes\partial_{z_1},
    \qquad
    f*g=m\circ\exp\!\left(\frac{\hbar_{\mathrm{W}}}{2}P\right)(f\otimes g),
  \]
  the normalized Moyal bracket
  \(\{f,g\}_{\hbar_{\mathrm{W}}}=\hbar_{\mathrm{W}}^{-1}(f*g-g*f)\) expands as
  \[
    \{f,g\}_{\hbar_{\mathrm{W}}}
      =
      \sum_{s\geq0}
      \frac{\hbar_{\mathrm{W}}^{2s}}{2^{2s}(2s+1)!}P^{2s+1}(f,g),
  \]
  with first quantum correction \(\hbar_{\mathrm{W}}^2P^3(f,g)/24\):
  \[
    \{f,g\}_{\hbar_{\mathrm{W}}}=\{f,g\}+\frac{\hbar_{\mathrm{W}}^2}{24}P^3(f,g)+O(\hbar_{\mathrm{W}}^4).
  \]
\end{lem}

\begin{proof}
  Even powers of \(P\) are symmetric and odd powers antisymmetric
  under \(f\leftrightarrow g\); the commutator therefore retains only
  the odd powers and doubles them:
  \[
    f*g-g*f=2\sum_{s\geq0}\frac{(\hbar_{\mathrm{W}}/2)^{2s+1}}{(2s+1)!}P^{2s+1}(f,g).
  \]
  Division by \(\hbar_{\mathrm{W}}\) gives the displayed series.  The \(s=1\)
  coefficient is \(1/(2^2\cdot3!)=1/24\); the all-order identity is
  Theorem~\ref{thm:app-radial-moyal-coefficients}, and exact rational
  coefficient checks through order eleven give the same sequence.
\end{proof}

\begin{lem}[Moyal coadjoint signs]
\label{lem:app-sign-moyal-coadjoint}
  Let \((n)_r=n(n-1)\cdots(n-r+1)\), with \((n)_0=1\), denote the
  falling factorial.  For \(f=z_1^pz_2^q\), the normalized Moyal
  coadjoint action on principal parts is
  \[
    f\cdot_{\hbar_{\mathrm{W}}}\rho_{a,b}
      =
      -\sum_{\substack{r\geq1\\ r\ \mathrm{odd}}}
      \frac{\hbar_{\mathrm{W}}^{r-1}}{2^{r-1}r!}
      C_r(p,q;a-p+r,b-q+r)\,
      \rho_{a-p+r,b-q+r},
  \]
  with negative-index targets and the scalar target \(\rho_{0,0}\)
  set to zero.  Here
  \[
    C_r(p,q;m,n)=
      \sum_{\alpha=0}^{r}(-1)^\alpha\binom r\alpha
      (p)_{r-\alpha}(q)_\alpha(m)_\alpha(n)_{r-\alpha}
  \]
  is the bidifferential coefficient of
  Theorem~\ref{thm:app-radial-moyal-coefficients}.  At \(r=1\) the
  formula recovers the residue coadjoint of
  Lemma~\ref{lem:app-sign-principal-part},
  \[
    z_1^pz_2^q\cdot\rho_{a,b}
      =-(pb-qa+p-q)\rho_{a-p+1,b-q+1}.
  \]
\end{lem}

\begin{proof}
  Apply the bidifferential \(P^r\) to
  \(z_1^pz_2^q\otimes z_1^mz_2^n\):
  Theorem~\ref{thm:app-radial-moyal-coefficients} gives
  \(P^r(f,g)=C_r(p,q;m,n)\,z_1^{p+m-r}z_2^{q+n-r}\).  By
  Lemma~\ref{lem:app-sign-moyal-coefficients} the normalized Moyal
  bracket keeps only odd \(r\), with coefficient
  \(\hbar_{\mathrm{W}}^{r-1}/(2^{r-1}r!)\).  The coadjoint convention
  \((f\cdot_{\hbar_{\mathrm{W}}}\rho)(g)=-\rho(\{f,g\}_{\hbar_{\mathrm{W}}})\) introduces the
  leading minus sign and determines the target principal part: pairing with
  \(\rho_{a,b}\) selects \((m,n)\) such that \(p+m-r=a\) and
  \(q+n-r=b\), i.e.\ \(m=a-p+r,\ n=b-q+r\).  At \(r=1\) the
  coefficient is
  \(C_1(p,q;a-p+1,b-q+1)=p(b-q+1)-q(a-p+1)=pb-qa+p-q\),
  reproducing Lemma~\ref{lem:app-sign-principal-part}.
\end{proof}

\begin{cor}[Universal residue-dual normalization]
\label{cor:app-sign-universal-dual-normalization}
  With the residue pairing
  \(\langle\rho_{a,b},z_1^mz_2^n\rangle=\delta_{a,m}\delta_{b,n}\) and
  the coadjoint convention
  \[
    (f\cdot\rho)(g)=-\rho(\{f,g\}),\qquad
    (f\cdot_{\hbar_{\mathrm{W}}}\rho)(g)=-\rho(\{f,g\}_{\hbar_{\mathrm{W}}}),
  \]
  the linear Hamiltonians act on \(\mathcal P\) by
  \[
    z_1\cdot\rho_{a,b}=-(b+1)\rho_{a,b+1},\qquad
    z_2\cdot\rho_{a,b}=(a+1)\rho_{a+1,b},
  \]
  while the polynomial descendants of
  Lemma~\ref{lem:app-sign-pbw-descendant} satisfy
  \[
    z_1\cdot\widetilde\Psi_{a,b}=b\,\widetilde\Psi_{a,b-1},\qquad
    z_2\cdot\widetilde\Psi_{a,b}=-a\,\widetilde\Psi_{a-1,b}.
  \]
  Diagonal rescaling
  \(\rho_{a,b}\mapsto\lambda_{a,b}\rho_{a,b}\),
  \(\widetilde\Psi_{a,b}\mapsto\mu_{a,b}\widetilde\Psi_{a,b}\) with
  \(\lambda_{a,b},\mu_{a,b}\in\C^\times\) preserves the opposite index
  directions of the two actions.  Adjustments of factorial normalizations or of the Moyal
  \(1/24\) coefficient act diagonally and leave the index shifts
  intact.
\end{cor}

\begin{proof}
  The two displayed actions and the residue pairing are restated from
  Lemma~\ref{lem:app-sign-principal-part} and
  Lemma~\ref{lem:app-sign-pbw-descendant} at
  \((p,q)\in\{(1,0),(0,1)\}\).  A diagonal rescaling sends
  \(z_1\cdot\rho_{a,b}\) to
  \(-(b+1)(\lambda_{a,b+1}/\lambda_{a,b})\rho_{a,b+1}\), preserving the
  raising target index, and sends
  \(z_1\cdot\widetilde\Psi_{a,b}\) to
  \(b(\mu_{a,b-1}/\mu_{a,b})\widetilde\Psi_{a,b-1}\), preserving the
  lowering target index.  Equality of the two actions would require the same
  target index for both, contradicting the directions
  \(b\mapsto b+1\) versus \(b\mapsto b-1\).  Higher Moyal corrections
  \(r\geq3\) shift indices by \(b\mapsto b-q+r\geq b+1\) for
  \((p,q)=(1,0)\) and by \(a\mapsto a-p+r\geq a+1\) for
  \((p,q)=(0,1)\); they raise as well, so the obstruction persists at
  every order in \(\hbar_{\mathrm{W}}\).  The finite-window homology calculation
  records this obstruction explicitly: writing \(\delta_{a,b}\) for the residue
  functional \(\rho_{a,b}\) under the Matlis identification of
  Lemma~\ref{lem:app-matlis-continuous-dual}, the third-order Moyal
  coadjoint gives
  \(\operatorname{ad}^*_{z_1^4}(\delta_{1,1})=-24\hbar_{\mathrm{W}}^2\delta_{0,4}\),
  while the polynomial one-\(\psi\) action sends \(\Psi_{1,1}\) to
  \(4\Psi_{4,0}\); even matched at the Capelli/Weyl shift \(\hbar_{\mathrm{W}} N\),
  the higher Moyal coadjoint terms have zero PBW descendant counterpart.
\end{proof}

\begin{lem}[Brane-line current source convention]
\label{lem:app-sign-current-source}
  For an interval \(I\subset\R\), the Hamiltonian source complex is
  \(\Omega^\bullet_c(I)\widehat\otimes\mathfrak h\) of
  Definition~\ref{def:app-factorization-current-source}.  A
  degree-zero source density is written \(a(t)dt\otimes f\), and the
  cotangent CE source coordinate maps to the boundary observable by
  \[
    u_{a(t)dt\otimes f}\longmapsto
    B_f(a)=\int_Ia(t)B_f(t)\,dt.
  \]
  The companion ghost coordinate \(c_{a\otimes f}\) maps to the PV
  coordinate \(\theta_{a\otimes f}\).  For
  \(B\in\mathcal D'_c(I)\), the multiplication of the distribution
  \(B\) by a smooth test function \(a\) is the
  \(C_c^\infty(I)\)-module operation \(aB\), defined by
  \((aB)(\varphi)=B(a\varphi)\).  The interval current brackets are
  \[
    \{B_f(a),\Theta_\rho(B)\}=\Theta_{f\cdot\rho}(aB),\qquad
    \{\Theta_\rho(B),\Theta_\sigma(C)\}=0.
  \]
\end{lem}

\begin{proof}
  The de Rham factor gives the compactly supported density on the
  brane line; the Hamiltonian current factor gives \(f\).  Pairing
  this source with the boundary field gives the integral defining
  \(B_f(a)\).  The first bracket is the semidirect-product bracket of
  Definition~\ref{def:app-factorization-current-dual} for
  \(\mathfrak h_I\ltimes\mathcal P_I[1]\): the action factors as the
  \(C_c^\infty(I)\)-module operation \(B\mapsto aB\) tensored with
  the Matlis coadjoint action of
  Lemma~\ref{lem:app-sign-principal-part}.  The second bracket
  vanishes because the shifted principal-part summand is an abelian
  odd ideal.
\end{proof}
